Prompts are edited all the time, yet most evaluations report results for a single prompt. We ask how much small, meaning-preserving edits to an instruction change task performance, whether the effect depends on task type, and whether combined edits compound. We define a Prompt Fragility Index (\pfi), the relative score drop caused by an edit. We apply six benign edit types (whitespace, capitalization, punctuation, typo, synonym, reorder), alone and combined, to the instructions of four tasks covering classification, multiple choice and generation. Our seeded editor builds each combined edit from the same local edits as its single-edit parts, so we measure the interaction term directly. Across 20,800 scored generations from two small instruction-tuned models (Qwen2.5-1.5B and Llama-3.2-1B), all three of our hypotheses are refuted. No edit type is uniformly harmful. For Qwen on 4-way topic classification, capitalization and punctuation edits cost 25 and 16 accuracy points, because they change how the label list is cased and delimited. Sentence reordering never costs more than 4.5 points. Generation is not the most fragile task type (mean \pfi 0.00--0.03, versus 0.10 for topic classification). Combined edits are super-additive only for generation (combined/additive ratio 1.29--1.61) and sub-additive elsewhere, because an item can break only once. Much apparent fragility is format fragility: with a longer output budget and a lenient parser, Llama's 28--48-point multiple-choice drops shrink to at most 10 points. Fragility follows what an edit touches, not what kind of edit it is.
